% ECONOMICS_PROBLEM: {"id": "C65-29", "title": "Nonconvex quadratic BSDE uniqueness: a necessary correction", "jel_code": "C65", "source_id": "OP-0566"}
% FIRST_POSTED: 2026-10-09
% ATTEMPT_STATUS: solved
% ENTRY_KIND: research_attempt
% !TeX program = pdflatex
% Self-contained: no external input or bibliography files required.
\documentclass[11pt,a4paper]{article}
\usepackage[T1]{fontenc}
\usepackage[utf8]{inputenc}
\usepackage{lmodern}
\usepackage[margin=24mm]{geometry}
\usepackage{amsmath,amssymb,amsthm,mathtools}
\usepackage{booktabs,longtable,array,enumitem,microtype,xurl,hyperref}
\hypersetup{colorlinks=true,linkcolor=blue,citecolor=blue,urlcolor=blue,pdftitle={Checked research attempts: nine economics and stochastic-analysis problems}}
\setlength{\parindent}{0pt}
\setlength{\parskip}{3pt}
\setlength{\emergencystretch}{3em}
\widowpenalty=10000\clubpenalty=10000\displaywidowpenalty=10000
\setlist[enumerate]{leftmargin=*,itemsep=2pt,topsep=3pt}
\newcommand{\R}{\mathbb R}
\newcommand{\N}{\mathbb N}
\newcommand{\E}{\mathbb E}
\newcommand{\Pp}{\mathbb P}
\newcommand{\Q}{\mathbb Q}
\newcommand{\F}{\mathcal F}
\newcommand{\ind}{\mathbf 1}
\newcommand{\cE}{\mathcal E}
\newcommand{\Law}{\operatorname{Law}}
\newcommand{\sgn}{\operatorname{sgn}}
\newcommand{\Var}{\operatorname{Var}}
\newcommand{\dist}{\operatorname{dist}}
\newcommand{\KL}{\operatorname{KL}}
\newcommand{\TV}{\operatorname{TV}}
\newcommand{\Lip}{\operatorname{Lip}}
\DeclareMathOperator*{\esssup}{ess\,sup}
\DeclareMathOperator*{\argmax}{arg\,max}
\newtheorem{theorem}{Theorem}
\newtheorem{lemma}[theorem]{Lemma}
\newtheorem{proposition}[theorem]{Proposition}
\newtheorem{corollary}[theorem]{Corollary}
\theoremstyle{definition}
\newtheorem{example}[theorem]{Example}
\newcommand{\report}[3]{\clearpage\setcounter{theorem}{0}\renewcommand{\thetheorem}{#1.\arabic{theorem}}\renewcommand{\theHtheorem}{#1.\arabic{theorem}}\section*{#1: #2}\addcontentsline{toc}{section}{#1: #2}\textbf{Outcome:} #3.\par\textbf{Tools:} web browsing available; Python computation available. Literature checked on 9 October 2026.}
\newcommand{\stage}[2]{\subsection*{#1) #2}}
\newcommand{\unresolved}{\textbf{LABEL: UNRESOLVED.}}
\newcommand{\disproof}{\textbf{LABEL: FULL SOLUTION. SUBLABEL: COUNTEREXAMPLE/DISPROOF.}}

\begin{document}
\section*{C65-29: Nonconvex BSDE uniqueness: literal versus repaired statement}
\textbf{Outcome:} FULL DISPROOF of the literal statement; repaired statement UNRESOLVED.\par\textbf{Tools:} web browsing available; Python computation available. Literature checked on 9 October 2026.
\subsection*{1) FORMAL RESTATEMENT}
\paragraph{Brownian conventions used in this report.}
Fix $T>0$ and an integer $d\ge1$. The probability space carries a $d$-dimensional Brownian motion $B$ with its completed, right-continuous natural filtration; $\F=\F_T$ and $\F_0$ is trivial modulo null sets. A scalar solution consists of a continuous adapted real process $Y$ and a predictable $\R^d$-valued process $Z$ such that
\[
 \int_0^T |Z_t|^2dt+\int_0^T|g(t,Y_t,Z_t)|dt<\infty\quad\Pp\text{-a.s.},
\]
and the displayed BSDE holds simultaneously for all times outside one null set. Equality of solutions means indistinguishability of $Y$ and $dt\otimes\Pp$-a.e. equality of $Z$. A process $X$ is of class (D) when $\{X_\tau:\tau\le T\text{ a stopping time}\}$ is uniformly integrable. Logs are natural, norms Euclidean, and $\sgn(0)=0$. At $T=0$ the assertions reduce to $Y_0=\xi$, with no nontrivial $Z$ component.

\paragraph{Precisely stated probabilistic tools.}
We use It\^o's formula and its Tanaka version for continuous semimartingales, Brownian martingale representation, and the following forms of change of measure. If $\theta$ is predictable and $\E\exp(\frac12\int_0^T|\theta|^2)<\infty$, Novikov's theorem says that
$D=\cE(\int\theta\,dB)$ is a uniformly integrable martingale; under $d\Q=D_Td\Pp$, $B-\int\theta dt$ is Brownian. Here $\cE(M)=\exp(M-\langle M\rangle/2)$. Also, for a continuous Brownian martingale, the BMO hypothesis
\[
 \esssup_{\tau\le T}\E\left[\left.\int_\tau^T|\theta_t|^2dt\right|\F_\tau\right]<\infty
\]
implies that $D$ is uniformly integrable and $D_T\in L^q$ for some $q>1$; this is the continuous-martingale BMO exponential theorem \cite{Kazamaki}. Each use below verifies the relevant integrability hypothesis. A local submartingale or supermartingale of class (D) is a true submartingale or supermartingale: apply the localized conditional inequalities and pass to the limit using uniform integrability. A local martingale of class (D) therefore equals the conditional expectation of its terminal value.

The literal assertion quantifies over progressively measurable continuous scalar generators, neither convex nor concave in $z$, with constants $C>0$ and $\delta\in[0,1]$ such that
\[
 |g(t,y,z)-g(t,y,z')|\le C(1+|z|^\delta+|z'|^\delta)|z-z'|.
\]
For every terminal value with $\E e^{\mu|\xi|}<\infty$ for every $\mu>0$, it asks for at most one solution of
\[
 Y_t=\xi+\int_t^Tg(s,Y_s,Z_s)ds-\int_t^TZ_s\,dB_s
\]
in the class in which $e^{\mu|Y|}$ is of class (D) for every $\mu>0$. At $\delta=0$ we interpret the growth bound by the constant $3C$, including at $z=0$.

The literal conditions put no uniqueness modulus on the $y$ dependence. The uploaded file already supplies a counterexample. The explicit repair in that same file is
\[
 |g(t,y,z)-g(t,y',z)|\le C|y-y'|,
 \qquad |g(t,0,0)|\le C.
\]
We keep this repair separate rather than pretending the displayed $z$ bound implies it. It is a sufficient conventional way to remove the elementary obstruction, not a claim that global Lipschitz continuity is the weakest possible repair.

\subsection*{2) QUICK LITERATURE/CONTEXT CHECK}
The source question is Problem 5.5 of Fan--Hu--Tang \cite{FHT}. The literal counterexample below is already present in the attachment, so it is verified rather than claimed as a new discovery. The nonconvex results in \cite{FHTnonconvex,WZF} impose further structural conditions; we do not substitute those conditions into the repaired target.

\subsection*{3) ATTACK PLAN}
\textbf{Proof routes.} Linearize two solutions, use a stochastic exponential to remove the $z$ difference, and test whether all exponential moments of $Y$ suffice for sampling under that probability. For bounded terminal values first obtain a deterministic bound and then BMO control.

\textbf{Disproof routes.} Set $Z=0$ and look for a non-Lipschitz scalar ODE; verify that an oscillatory $z$ term makes the generator neither convex nor concave; after adding the repair, test unbounded-$Z$ cases rather than reusing the excluded ODE example.

\textbf{Landscape and tools.} This is stochastic uniqueness with a deterministic obstruction. The useful tools are separation of variables (for delayed ODE solutions), radial derivative tests (for nonconvexity), divided differences (for exact linearization), Girsanov's theorem (for uniqueness of the linearized equation), H\"older's inequality (for class-(D) transfer), Tanaka/exponential transforms (for bounded terminal estimates), and BMO exponentials (for unbounded but controlled integrands).

\subsection*{4) WORK}
\begin{theorem}[Literal statement is false]
For every $T>0$ and $d\ge1$, let
\[
 g(y,z)=\sqrt{|y|}+|z|^2+3\sin|z|,\qquad \xi=0.
\]
For every $a\in[0,T]$, the pair
\[
 Z^a_t=0,\qquad Y^a_t=\tfrac14((a-t)^+)^2
\]
is an admissible solution. Distinct $a$ give distinct solutions.
\end{theorem}
\begin{proof}
Write $r=|z|$ and $r'=|z'|$. The reverse triangle inequality and the $1$-Lipschitz property of sine give
\[
 |r^2-(r')^2+3\sin r-3\sin r'|
 \le(r+r'+3)|z-z'|
 \le3(1+r+r')|z-z'|.
\]
Thus the required inequality holds with $C=3$, $\delta=1$. On a positive coordinate ray, the second derivative of the $z$ part is $2-3\sin r$, equal to $-1$ at $r=\pi/2$ and $5$ at $r=3\pi/2$. A twice differentiable convex function has nonnegative second derivative and a concave one nonpositive second derivative on every open interval. Restriction to this ray therefore excludes both convexity and concavity. The generator is continuous and deterministic, hence progressively measurable.

Since $g(Y^a_s,0)=\sqrt{Y^a_s}=(a-s)^+/2$, direct integration gives
\[
 \int_t^Tg(Y^a_s,0)ds=\int_t^a\frac{a-s}{2}ds
 =\frac{(a-t)^2}{4}\quad(t<a),
\]
and both sides are zero when $t\ge a$. This is exactly the BSDE with zero terminal value and zero martingale term. Both required time integrals are finite. Moreover
$0\le Y^a_t\le T^2/4$, so for every $\mu>0$ the entire stopping family $e^{\mu|Y^a_\tau|}$ is bounded by $e^{\mu T^2/4}$ and is uniformly integrable. At $t=0$, $Y^a_0=a^2/4$, so distinct $a$ give distinct solutions. Finally $\sqrt{|y|}\le1+|y|$, but
$|\sqrt y-0|/|y-0|=y^{-1/2}\to\infty$ as $y\downarrow0$. The example has at most linear $y$ growth and fails precisely the added Lipschitz repair.
\end{proof}

The remaining results concern only the repaired question.

\begin{lemma}[Transfer under an integrable density]
If $D_T\in L^q(\Pp)$ for some $q>1$ and
$\sup_\tau\E|X_\tau|^p<\infty$ for every finite $p$, then $X$ is of class (D) under $d\Q=D_Td\Pp$. If the density integrand is uniformly bounded, the same conclusion holds whenever the stopping family has a uniformly bounded $p$th moment for one $p>1$.
\end{lemma}
\begin{proof}
In the first case choose $r>1$ and $p=rq/(q-1)$. H\"older gives
\[
 \sup_\tau\E_{\Q}|X_\tau|^r
 \le\|D_T\|_{L^q}\,
       \sup_\tau\bigl(\E|X_\tau|^p\bigr)^{(q-1)/q}<\infty.
\]
Uniformly bounded $r$th moments imply uniform integrability. For a bounded integrand $|b|\le M$, Novikov applies to every multiple of $b$ and
\[
 \E D_T^q\le \exp\bigl(q(q-1)M^2T/2\bigr),\qquad q>1,
\]
because $D_T^q=\cE(q\int b\,dB)_T
\exp(\frac12q(q-1)\int|b|^2)$.
For a given $p>1$, choose $1<r<p$ and then $q=p/(p-r)$ in the same H\"older inequality. This proves the second assertion. The all-exponential class in the problem has all uniformly bounded stopping moments, since for every finite $p$ there is a constant $K_p$ with $x^p\le K_pe^x$ for $x\ge0$.
\end{proof}

\begin{proposition}[A precise linearization criterion]
For two solutions to the repaired equation, set $\Delta Y=Y-\widetilde Y$ and $\Delta Z=Z-\widetilde Z$. Define
\[
 a_t=\frac{g(t,Y_t,Z_t)-g(t,\widetilde Y_t,Z_t)}{Y_t-\widetilde Y_t}
\]
where the denominator is nonzero, and zero otherwise, and define
\[
 b_t=\frac{g(t,\widetilde Y_t,Z_t)-g(t,\widetilde Y_t,\widetilde Z_t)}
 {|\Delta Z_t|^2}\,\Delta Z_t
\]
where $\Delta Z_t\ne0$, and zero otherwise. If $\cE(\int b\,dB)$ is a true density with terminal value in some $L^q$, $q>1$, then the two solutions coincide.
\end{proposition}
\begin{proof}
Measurability follows from the measurable generator and the explicit zero-denominator convention. The hypotheses imply $|a|\le C$ and
$|b|\le C(1+|Z|^\delta+|\widetilde Z|^\delta)$, so $\int|b|^2<\infty$ a.s. by $\delta\le1$ and finite $T$. The difference equation is exactly
$d\Delta Y=-(a\Delta Y+b\cdot\Delta Z)dt+\Delta Z\,dB$.
Under the stated probability $\Q$, the process
$L_t=\exp(\int_0^ta_sds)\Delta Y_t$ is a local martingale. The transfer lemma and the deterministic bound on the multiplier make $L$ of class (D). Its terminal value is zero, so $L_t=\E_\Q[L_T\mid\F_t]=0$. Equality of $Z$ follows from quadratic variation as before.
\end{proof}

\begin{corollary}[Solved subcases of the repair]
Uniqueness holds for $\delta=0$. It also holds if both solutions' $Z$ integrands are uniformly bounded, or if the linearization integrand $b$ satisfies the BMO hypothesis.
\end{corollary}
\begin{proof}
For $\delta=0$ the bound is $|b|\le3C$. If both $Z$ are bounded the same reasoning gives a finite deterministic bound on $b$. The transfer lemma verifies the density condition. In the BMO case the stated continuous-martingale BMO theorem supplies it.
\end{proof}

\begin{theorem}[Bounded terminal data for every $0\le\delta\le1$]
Under the repaired assumptions and $|\xi|\le K$, there is at most one solution in the all-exponential class-(D) class.
\end{theorem}
\begin{proof}
There are deterministic constants $L\ge1$, $\gamma>0$ such that
\[
 |g(t,y,z)|\le L(1+|y|)+\frac\gamma2|z|^2.
\]
Indeed the hypotheses bound the left side by
$C+C|y|+C(|z|+|z|^{1+\delta})$, and both $r$ and $r^{1+\delta}$ are bounded by a constant times $1+r^2$ for $0\le\delta\le1$.

Tanaka's formula gives
$d|Y|\ge-[L(1+|Y|)+\gamma|Z|^2/2]dt+\sgn(Y)Z\,dB+dL^0(Y)$,
where the local-time term is nonnegative (a choice of normalization only rescales $L^0$). For a continuous semimartingale the occupation-density formula gives
$\int_0^T\ind_{\{Y_t=0\}}|Z_t|^2dt=0$; thus the martingale part of $|Y|$ has quadratic variation $\int|Z|^2dt$.
Apply It\^o to
\[
 H_t=\exp\{\gamma[e^{Lt}|Y_t|+(e^{Lt}-1)]\}.
\]
The terms involving $L(1+|Y|)$ cancel, and the remaining quadratic drift is bounded below by
$\frac12\gamma^2(e^{2Lt}-e^{Lt})H_t|Z_t|^2dt\ge0$.
The local-time contribution is also nonnegative. This makes $H$ a local submartingale. Its stopping family is dominated by a constant times $\exp(\gamma e^{LT}|Y_\tau|)$, so it is of class (D). Consequently
\[
 H_t\le\E[H_T\mid\F_t]
 \le\exp\{\gamma[e^{LT}K+(e^{LT}-1)]\},
\]
which implies
$|Y_t|\le e^{L(T-t)}(K+1)-1\le K':=e^{LT}(K+1)-1$.

Next choose $a>\gamma$. The drift in It\^o's formula for $e^{aY}$ is bounded below by
\[
 e^{aY}\left[-aL(1+K')+\frac{a(a-\gamma)}2|Z|^2\right]dt.
\]
Integrate from any stopping time $\tau$ to localized times increasing to $T$, take conditional expectations, and use $e^{-aK'}\le e^{aY}\le e^{aK'}$. Conditional Fatou for the nonnegative $Z$ term gives
\[
 \E\left[\left.\int_\tau^T|Z_t|^2dt\right|\F_\tau\right]
 \le\frac{2e^{2aK'}[2+aL(1+K')T]}{a(a-\gamma)}.
\]
This finite deterministic upper bound proves BMO. It holds for both solutions. Since $r^{2\delta}\le1+r^2$, the previously defined $b$ has $|b|^2\le K''(1+|Z|^2+|\widetilde Z|^2)$ for a finite constant $K''$. It too is BMO. The linearization criterion now proves uniqueness.
\end{proof}

\subsection*{5) VERIFICATION}
The counterexample includes $a=0$ and $a=T$ and has no stochastic integrability issue. At the join $t=a$, its derivative is zero on both sides. The script directly compares the integral and solution at four choices of $a$ and 21 time points; maximum discrepancy is $5.6\cdot10^{-17}$. The radial second derivatives used above are exactly $-1$ and $5$.

For the repaired question, the bound on $b$ is local in $Z$, not a deterministic global Lipschitz constant when $\delta>0$. All exponential moments of $Y$ do not, by the reasoning supplied here, automatically imply BMO of $Z$ for unbounded terminal data. The bounded-terminal proof uses all-exponential class (D) at the precise transform exponent $\gamma e^{LT}$; dropping that hypothesis would invalidate that step. The BMO theorem is used only for continuous martingales, not for a general discontinuous process.

\subsection*{6) FINAL}
\textbf{Literal statement: LABEL: FULL SOLUTION. SUBLABEL: COUNTEREXAMPLE/DISPROOF.}
The generator and continuum of solutions in the first theorem provide a fully verified disproof. This counterexample was already supplied in the uploaded file.

\textbf{Repaired statement: LABEL: UNRESOLVED.}
The strongest proved subclasses are $\delta=0$, bounded terminal data for all $\delta\in[0,1]$, and the stated bounded-integrand/BMO comparison criteria.

\textbf{Exact first gap.} For arbitrary repaired data and two all-exponential admissible solutions, establish a valid change of measure and terminal sampling for
$e^{\int a}\Delta Y$ with the unbounded divided-difference integrand $b$. A sufficient unproved lemma is that this exponential density is a true martingale with some $L^q$, $q>1$, terminal moment; an alternative proof could avoid that sufficient lemma.

\textbf{Next three moves.} (1) Derive quantitative exponential-energy estimates for $\int|Z|^{2\delta}$ from the all-exponential solution class, with constants sufficient for sampling. (2) Separate $0<\delta<1$ from $\delta=1$ and test stopped linearizations at the exact required moment order. (3) Search generators whose derivative oscillates at order $|z|$ together with unbounded terminal tails, verifying both all-exponential class-(D) requirements before calling any construction a counterexample.

\textbf{Likely minimal counterexample to the repair, if any.} It must have unbounded terminal data and escape the bounded-density/BMO cases. A $Z=0$ delayed ODE construction is excluded by global Lipschitz continuity in $y$ and cannot answer the repaired question.

\clearpage
\begin{thebibliography}{99}
\bibitem{Kazamaki}
Norihiko Kazamaki. \emph{Continuous Exponential Martingales and BMO}. Lecture Notes in Mathematics 1579, Springer, 1994. DOI: 10.1007/BFb0073585. Continuous BMO stochastic exponentials and reverse H\"older inequalities. \url{https://doi.org/10.1007/BFb0073585}.

\bibitem{FHT}
Shengjun Fan, Ying Hu and Shanjian Tang. \emph{1D nonlinear backward stochastic differential equations: a unified theory and applications}. arXiv:2309.06233v2, 11 February 2026. In particular Problems 5.1--5.5, PDF pp.\ 32--33; for C65-45 see (A1), (4.13)--(4.14), PDF p.\ 31. \url{https://arxiv.org/pdf/2309.06233v2}.

\bibitem{FHTnonconvex}
Shengjun Fan, Ying Hu and Shanjian Tang. \emph{On the uniqueness of solutions to quadratic BSDEs with non-convex generators and unbounded terminal conditions}. Comptes Rendus Math\'ematique 358(2) (2020), 227--235. DOI: 10.5802/crmath.40. \url{https://www.numdam.org/articles/10.5802/crmath.40/}.

\bibitem{WZF}
Yan Wang, Yaqi Zhang and Shengjun Fan. \emph{On the uniqueness of solutions to quadratic BSDEs with non-convex generators and unbounded terminal conditions: the certain exponential moment case}. arXiv:2409.13463v2, 3 August 2026. \url{https://arxiv.org/abs/2409.13463v2}.

\end{thebibliography}
\end{document}
